\chapter{Graphs, terms and bounded reduction}\label{ch:8}
\section*{The working problem}
Separate graph storage from term semantics, then choose a reduction strategy.

DirectedGraph provides node and edge identity. Terms provide ordered function arguments and variables. A term rewrite system uses matching, not arbitrary evaluation. Bounded normalization reports whether it found a normal form or exhausted its budget.

\section*{Graph algorithms and shared reduction}
DirectedGraph supports breadth-first and depth-first traversal, topological sorting, strongly connected components and DOT output. A multigraph can contain parallel edges, so algorithms and tests should not silently assume one edge per source-target pair. IDs can become invalid after removal; graph storage uses optional entries rather than recycling every hole into an unrelated node. Keep validity checks at plugin and external-data boundaries.

\section*{In this chapter}
\begin{enumerate}
\item Graph identity and edges
\item Dependency order
\item Matching first-order terms
\item Normalization with an explicit budget
\item Shared term graphs
\end{enumerate}
Read the workshops in order, then use the chapter review to connect the small examples to an application boundary. Each example is independently runnable; it does not depend on variables introduced in an earlier listing.

\clearpage
\workshop{8.1}{Graph identity and edges}{C++20}
A directed graph represents relationships with stable IDs rather than child-value copies. Use edge payloads when relationships themselves have roles, such as operand positions or dependency labels. Keep node IDs local to the graph that created them.

\lstinputlisting[language=C++,caption={Example ex08\_01: C++}]{examples/cpp/ex08_01.inc}
\section*{Execution, step by step}
The edge stores a relationship between IDs. Neither node becomes a nested value inside the other.

\section*{Read the implementation}
The relevant source entry is DirectedGraph in vendor/DSLtk.hpp, around line 11527. Follow the public call into its result construction or dispatch branch. Identify the input ownership, the success condition and the exact failure path before changing this example.

The assertions are the acceptance contract for this workshop. Run the complete program, then change one input at a time. When a change fails, first identify whether the failure belongs to the native build, structural recognition, language checking, execution or the domain result.

\clearpage
\workshop{8.1}{Graph identity and edges}{D}
This D facility is a reference type, so construction uses new. Copying the reference aliases the same mutable object; it does not perform the AST-style deep copy demonstrated earlier.

\lstinputlisting[language=D,caption={Example ex08\_01: D}]{examples/d/ex08_01.inc}
\section*{Failure and design boundary}
Removing a node invalidates its ID for future access. A small integer from a different graph is not a valid substitute.

\section*{Try it yourself}
Add a second outgoing edge and compare incoming, outgoing and successor lists.

\section*{Review the result}
Keep the original accepted fixture and add the new boundary fixture beside it. Compare observable values and diagnostic categories across the two hosts. Do not label the D run as passed until a D compiler has executed it; the supplied edition records source review separately from runtime evidence.

\clearpage
\workshop{8.2}{Dependency order}{C++20}
Topological sorting is useful for acyclic dependency graphs. It answers a scheduling question, whereas breadth-first and depth-first traversal answer reachability questions. Do not substitute one algorithm merely because all of them return a vector of IDs.

\lstinputlisting[language=C++,caption={Example ex08\_02: C++}]{examples/cpp/ex08_02.inc}
\section*{Execution, step by step}
The producer appears before its dependent. With several independent nodes, more than one valid order may exist.

\section*{Read the implementation}
The relevant source entry is DirectedGraph in vendor/DSLtk.hpp, around line 11527. Follow the public call into its result construction or dispatch branch. Identify the input ownership, the success condition and the exact failure path before changing this example.

The assertions are the acceptance contract for this workshop. Run the complete program, then change one input at a time. When a change fails, first identify whether the failure belongs to the native build, structural recognition, language checking, execution or the domain result.

\clearpage
\workshop{8.2}{Dependency order}{D}
This D facility is a reference type, so construction uses new. Copying the reference aliases the same mutable object; it does not perform the AST-style deep copy demonstrated earlier.

\lstinputlisting[language=D,caption={Example ex08\_02: D}]{examples/d/ex08_02.inc}
\section*{Failure and design boundary}
Test cycles explicitly and inspect the implementation contract. A dependency scheduler must not execute a partial order as though every node was scheduled.

\section*{Try it yourself}
Create a diamond dependency and verify every edge respects the resulting order.

\section*{Review the result}
Keep the original accepted fixture and add the new boundary fixture beside it. Compare observable values and diagnostic categories across the two hosts. Do not label the D run as passed until a D compiler has executed it; the supplied edition records source review separately from runtime evidence.

\clearpage
\workshop{8.3}{Matching first-order terms}{C++20}
Term matching binds variables in a pattern against a fixed subject. Repeated occurrences of one variable must agree. This is a structural operation and should not be confused with solving an arithmetic equation.

\lstinputlisting[language=C++,caption={Example ex08\_03: C++}]{examples/cpp/ex08_03.inc}
\section*{Execution, step by step}
The first argument binds x. The second occurrence checks that the subject agrees with that binding.

\section*{Read the implementation}
The relevant source entry is match\_term in vendor/DSLtk.hpp, around line 9604. Follow the public call into its result construction or dispatch branch. Identify the input ownership, the success condition and the exact failure path before changing this example.

The assertions are the acceptance contract for this workshop. Run the complete program, then change one input at a time. When a change fails, first identify whether the failure belongs to the native build, structural recognition, language checking, execution or the domain result.

\clearpage
\workshop{8.3}{Matching first-order terms}{D}
The C++ optional test becomes D Nullable.isNull, with the sense reversed. Test the outer parse result before accessing an inner optional payload; absence and parse failure can occupy different layers.

\lstinputlisting[language=D,caption={Example ex08\_03: D}]{examples/d/ex08_03.inc}
\section*{Failure and design boundary}
Matching is directional. Unification allows variables on both sides and needs an occurs check to prevent infinite substitutions.

\section*{Try it yourself}
Compare match\_term and unify\_terms on a pair of terms that both contain variables.

\section*{Review the result}
Keep the original accepted fixture and add the new boundary fixture beside it. Compare observable values and diagnostic categories across the two hosts. Do not label the D run as passed until a D compiler has executed it; the supplied edition records source review separately from runtime evidence.

\clearpage
\workshop{8.4}{Normalization with an explicit budget}{C++20}
A rewrite system repeatedly applies its rules according to a selected traversal strategy. The step budget makes nontermination observable. A returned term is always available, but callers must inspect the stop reason before claiming that it is normalized.

\lstinputlisting[language=C++,caption={Example ex08\_04: C++}]{examples/cpp/ex08_04.inc}
\section*{Execution, step by step}
The only rule removes the add node. A further search finds no redex, establishing a normal form for this rule set.

\section*{Read the implementation}
The relevant source entry is TermRewriteSystem in vendor/DSLtk.hpp, around line 11920. Follow the public call into its result construction or dispatch branch. Identify the input ownership, the success condition and the exact failure path before changing this example.

The assertions are the acceptance contract for this workshop. Run the complete program, then change one input at a time. When a change fails, first identify whether the failure belongs to the native build, structural recognition, language checking, execution or the domain result.

\clearpage
\workshop{8.4}{Normalization with an explicit budget}{D}
This D facility is a reference type, so construction uses new. Copying the reference aliases the same mutable object; it does not perform the AST-style deep copy demonstrated earlier.

\lstinputlisting[language=D,caption={Example ex08\_04: D}]{examples/d/ex08_04.inc}
\section*{Failure and design boundary}
Normal form is relative to the provided rules. It does not mean globally optimal, semantically checked or canonical across different rule orders.

\section*{Try it yourself}
Add opposing rules and use a small budget. Confirm StepLimit is reported rather than normal form.

\section*{Review the result}
Keep the original accepted fixture and add the new boundary fixture beside it. Compare observable values and diagnostic categories across the two hosts. Do not label the D run as passed until a D compiler has executed it; the supplied edition records source review separately from runtime evidence.

\clearpage
\workshop{8.5}{Shared term graphs}{C++20}
A TermGraph can intern repeated terms so identical subexpressions share nodes. This changes storage and rewrite behavior: replacing a shared redex may affect several incoming references. Treat sharing as part of the execution model rather than a free optimization.

\lstinputlisting[language=C++,caption={Example ex08\_05: C++}]{examples/cpp/ex08_05.inc}
\section*{Execution, step by step}
Interning stores only the atom a, the application f(a), and the pair root. Both pair edges refer to the same f(a) node, so the count is three.

\section*{Read the implementation}
The relevant source entry is TermGraph in vendor/DSLtk.hpp, around line 11983. Follow the public call into its result construction or dispatch branch. Identify the input ownership, the success condition and the exact failure path before changing this example.

The assertions are the acceptance contract for this workshop. Run the complete program, then change one input at a time. When a change fails, first identify whether the failure belongs to the native build, structural recognition, language checking, execution or the domain result.

\clearpage
\workshop{8.5}{Shared term graphs}{D}
This D facility is a reference type, so construction uses new. Copying the reference aliases the same mutable object; it does not perform the AST-style deep copy demonstrated earlier.

\lstinputlisting[language=D,caption={Example ex08\_05: D}]{examples/d/ex08_05.inc}
\section*{Failure and design boundary}
Sharing must preserve pure term meaning. Duplicated effectful computations cannot generally be merged without changing observable behavior.

\section*{Try it yourself}
Repeat construction with sharing disabled and compare node counts and reconstructed terms.

\section*{Review the result}
Keep the original accepted fixture and add the new boundary fixture beside it. Compare observable values and diagnostic categories across the two hosts. Do not label the D run as passed until a D compiler has executed it; the supplied edition records source review separately from runtime evidence.

\clearpage
\section*{Chapter review: Graph algorithms and shared reduction}
Term unification differs from matching because variables may occur on both sides. The occurs check prevents binding a variable to a term containing itself. TermGraph reduction redirects incoming edges to a replacement and therefore implements shared reduction. Garbage collection takes explicit roots and removes unreachable nodes. Reconstructing a shared graph as a tree can duplicate work or storage, so use to\_term for inspection deliberately rather than in every hot-path operation.

\section*{Questions to answer in your implementation}
\begin{enumerate}
\item Add a second outgoing edge and compare incoming, outgoing and successor lists.
\item Create a diamond dependency and verify every edge respects the resulting order.
\item Compare match\_term and unify\_terms on a pair of terms that both contain variables.
\item Add opposing rules and use a small budget. Confirm StepLimit is reported rather than normal form.
\item Repeat construction with sharing disabled and compare node counts and reconstructed terms.
\end{enumerate}
\section*{A practical completion test}
Choose one workshop and replace its fixed input with an application boundary: a command-line argument, a file, a configuration field or a typed callback. Write down the accepted input shape and the error outcome. Keep parsing and validation separate from publishing state. Preserve the smallest failing input when a regression is found.

